\appendix

\chapter{团队分工}
\label{ch:team}

本大作业由4名成员协作完成。项目采用"技术负责人主导开发+组员分工交付"的模式，技术核心（数据采集、模型训练、技术报告撰写）由技术负责人完成，组员负责报告校对、演示视频、答辩PPT与团队报告等交付物。各成员具体分工与贡献如表\ref{tab:team-division}所示。

\begin{table}[H]
    \centering
    \caption{团队分工与贡献表}
    \label{tab:team-division}
    \begin{tabularx}{\textwidth}{clXc}
        \toprule
        成员 & 角色 & 主要负责内容 & 贡献比 \\
        \midrule
        成员1 & 技术负责人 & 数据采集（Pexels视频+ESC-50音频）、MediaPipe Tasks API迁移、数据预处理重构（900样本）、模型训练（Full 100\%/Lite 94.44\%）、消融实验（A/B/C三组）、端到端验证、源代码整理、\textbf{技术报告全部8章撰写+12篇参考文献} & 40\% \\
        成员2 & 报告校对与文献 & 报告全文校对润色、文献深度补充（2--4篇）、引言与相关工作创新点论述增强、LaTeX排版修正 & 20\% \\
        成员3 & 演示与可视化 & 演示视频制作（3--5分钟）、训练曲线图、消融对比柱状图、模态重要性排名图、混淆矩阵可视化、App截图 & 20\% \\
        成员4 & 答辩与展示 & 答辩PPT（5页）、团队贡献报告、答辩演讲稿、Q\&A准备、目录清理与最终打包 & 20\% \\
        \bottomrule
    \end{tabularx}
\end{table}

\section{已完成工作说明}

截至报告撰写日，以下工作已由技术负责人（成员1）完成：

\begin{enumerate}
    \item \textbf{数据层}：采集Pexels视频112段（daily 43/struggle 30/falls 39）+ ESC-50音频11类，构建900样本多模态数据集（训练720/验证90/测试90）。
    \item \textbf{特征层}：将视频特征从零向量升级为MediaPipe PoseLandmarker真实骨架特征（33关键点$\to$768维），非零率从7.6\%提升至100\%。
    \item \textbf{模型层}：训练Full模型（Cross-Modal Attention，100\%测试准确率）和Lite基线（MLP拼接，94.44\%），参数量分别为8.51M和3.74M。
    \item \textbf{实验层}：完成A组（融合策略）、B组（单模态屏蔽）、C组（模态重要性）三组消融实验，验证四模态均参与决策。
    \item \textbf{应用层}：Streamlit Web应用端到端验证6/6正确，置信度91\%$\to$95\%$\to$97\%。
    \item \textbf{报告层}：基于标准LaTeX模板完成8章技术报告（50页）+ 12篇参考文献，编译通过无警告。
\end{enumerate}

\section{剩余交付物}

以下工作由组员分工完成：

\begin{table}[H]
    \centering
    \caption{剩余交付物与负责人}
    \label{tab:remaining}
    \begin{tabularx}{\textwidth}{clXc}
        \toprule
        \# & 负责人 & 交付物 & 对应分值 \\
        \midrule
        1 & 成员1 & 源代码整理+复现脚本+README & 5 \\
        2 & 成员2 & 报告校对润色+文献补充 & 校对20 \\
        3 & 成员3 & 演示视频（3--5分钟） & 5 \\
        4 & 成员3 & 实验图表集（PPT/视频用） & 支撑30 \\
        5 & 成员4 & 答辩PPT（5页） & 5 \\
        6 & 成员4 & 团队贡献报告 & 5 \\
        7 & 成员4 & 答辩讲稿+Q\&A & --- \\
        \bottomrule
    \end{tabularx}
\end{table}

所有代码、模型、数据均在项目仓库中统一管理，确保可追溯与可复现。


\chapter{环境配置与复现指南}
\label{ch:reproducibility}

\section{硬件与软件环境}

\begin{table}[H]
    \centering
    \caption{运行环境}
    \label{tab:env}
    \begin{tabular}{ll}
        \toprule
        项目 & 配置 \\
        \midrule
        操作系统 & Windows 11 \\
        Python版本 & 3.10+ \\
        深度学习框架 & PyTorch 2.0+ \\
        视频分析 & MediaPipe Tasks API 0.10.35 \\
        音频分析 & librosa 0.10+ \\
        Web框架 & Streamlit 1.28+ \\
        训练设备 & CPU（Intel i7），未使用GPU \\
        \bottomrule
    \end{tabular}
\end{table}

\section{安装依赖}

\begin{lstlisting}[language=bash,caption={安装依赖}]
# 克隆项目后进入目录
cd smart_elderly_care_v2
pip install -r requirements.txt

# 应用端依赖
cd ../smart_elderly_care_app
pip install -r requirements.txt
\end{lstlisting}

\section{数据准备}

\begin{lstlisting}[language=bash,caption={下载ESC-50与Pexels视频数据集}]
cd smart_elderly_care_v2
python data/download_datasets.py
# 自动下载 ESC-50 (约600MB) 和 Pexels 视频 (约150MB)
\end{lstlisting}

\begin{lstlisting}[language=bash,caption={数据预处理（特征提取与数据集构建）}]
python data/preprocess_real.py
# 输出至 data/processed/，包含 train/val/test 三个子集
\end{lstlisting}

\section{模型训练}

\begin{lstlisting}[language=bash,caption={训练Full模型（Cross-Modal Attention）}]
python train.py --model full
# 训练日志输出至控制台，最优模型保存至 checkpoints/best_model.pt
\end{lstlisting}

\begin{lstlisting}[language=bash,caption={训练Lite模型（消融对照）}]
python train.py --model lite
# 模型保存至 checkpoints/lite/best_model.pt
\end{lstlisting}

\section{评估与消融}

\begin{lstlisting}[language=bash,caption={模型评估}]
python evaluate.py
# 输出准确率、精确率、召回率、F1、混淆矩阵
# 结果保存至 eval_results/
\end{lstlisting}

\begin{lstlisting}[language=bash,caption={消融实验}]
python ablation.py
# 运行A/B/C三组消融实验
# 结果保存至 ablation_results/
\end{lstlisting}

\section{应用运行}

\begin{lstlisting}[language=bash,caption={启动Streamlit应用}]
cd smart_elderly_care_app

# 首次运行需生成默认特征文件
python _regen_defaults.py

# 启动应用
streamlit run app.py
# 浏览器访问 http://localhost:8501
\end{lstlisting}

\begin{lstlisting}[language=bash,caption={端到端推理测试}]
python test_inference.py
# 使用真实ESC-50音频进行端到端推理验证
\end{lstlisting}

\section{项目文件结构}

\begin{lstlisting}[language=bash,caption={项目目录结构}]
multimodal/
|-- report/                          # LaTeX技术报告
|   |-- main.tex
|   |-- chapters/                    # 各章节
|   `-- bib/references.bib           # 参考文献
|-- smart_elderly_care_v2/           # 模型训练与评估
|   |-- config.py                    # 配置
|   |-- train.py                     # 训练脚本
|   |-- evaluate.py                  # 评估脚本
|   |-- ablation.py                  # 消融实验
|   |-- inference.py                 # 推理接口
|   |-- models/                      # 模型定义
|   |   |-- fusion_net.py            # 融合网络
|   |   |-- encoders.py              # 单模态编码器
|   |   |-- unified_feature_extractor.py  # 统一特征提取器
|   |   `-- dataset.py               # 数据集类
|   |-- data/                        # 数据处理
|   |   |-- download_datasets.py     # 数据集下载
|   |   |-- preprocess_real.py       # 数据预处理
|   |   `-- processed/               # 处理后数据
|   |-- checkpoints/                 # 模型权重
|   |-- eval_results/                # 评估结果
|   `-- ablation_results/            # 消融结果
|-- smart_elderly_care_app/          # Web应用
|   |-- app.py                       # Streamlit主程序
|   |-- inference/                   # 推理模块
|   |   |-- predictor.py             # 风险预测器
|   |   |-- feature_extractor.py     # 特征提取代理
|   |   `-- missing_handler.py       # 缺失数据处理
|   |-- pretrained_models/           # 部署模型
|   `-- defaults/                    # 默认特征文件
`-- plans/                           # 规划文档
\end{lstlisting}


\chapter{关键实验数据}
\label{ch:data}

\section{训练历史}

表\ref{tab:training-history-full}列出了Full模型完整训练过程（共15个epoch，第5个epoch达到最优后早停）。

\begin{table}[H]
    \centering
    \caption{Full模型训练历史}
    \label{tab:training-history-full}
    \begin{tabular}{ccccc}
        \toprule
        Epoch & 训练损失 & 训练准确率 & 验证损失 & 验证准确率 \\
        \midrule
        1 & 0.9537 & 56.67\% & 0.7234 & 73.33\% \\
        3 & 0.5821 & 78.89\% & 0.3845 & 88.89\% \\
        5 & 0.2943 & 93.33\% & 0.1572 & 100.00\% \\
        10 & 0.1024 & 98.89\% & 0.0613 & 100.00\% \\
        15 & 0.0387 & 100.00\% & 0.0291 & 100.00\% \\
        \bottomrule
    \end{tabular}
\end{table}

\section{测试集混淆矩阵}

Full模型在测试集（90样本）上的混淆矩阵如表\ref{tab:cm-test}所示，三个类别均达到100\%识别率。

\begin{table}[H]
    \centering
    \caption{测试集混淆矩阵（Full模型）}
    \label{tab:cm-test}
    \begin{tabular}{c|ccc|c}
        \toprule
        \multirow{2}{*}{真实类别} & \multicolumn{3}{c|}{预测类别} & \multirow{2}{*}{合计} \\
         & 低风险 & 中风险 & 高风险 & \\
        \midrule
        低风险 & 29 & 0 & 0 & 29 \\
        中风险 & 0 & 28 & 0 & 28 \\
        高风险 & 0 & 0 & 33 & 33 \\
        \midrule
        合计 & 29 & 28 & 33 & 90 \\
        \bottomrule
    \end{tabular}
\end{table}

\section{端到端推理测试结果}

表\ref{tab:e2e-test}列出了应用层端到端推理测试的详细结果，6组真实ESC-50音频样本全部预测正确。

\begin{table}[H]
    \centering
    \caption{端到端推理测试结果}
    \label{tab:e2e-test}
    \begin{tabular}{lllcc}
        \toprule
        编号 & 音频描述 & 预期风险 & 预测风险 & 置信度 \\
        \midrule
        1 & 狗叫（环境声） & 低风险 & 低风险 & 91.2\% \\
        2 & 雨声（自然声） & 低风险 & 低风险 & 90.8\% \\
        3 & 咳嗽（人声） & 中风险 & 中风险 & 94.7\% \\
        4 & 婴儿哭声（人声） & 中风险 & 中风险 & 95.3\% \\
        5 & 尖叫（人声） & 高风险 & 高风险 & 96.8\% \\
        6 & 玻璃碎裂（撞击声） & 高风险 & 高风险 & 97.4\% \\
        \bottomrule
    \end{tabular}
\end{table}

置信度随风险等级递增（91\%$\to$95\%$\to$97\%），符合高风险事件信号更显著的直觉，间接验证了模型学习到的决策边界具有临床合理性。
